
\documentclass[aspectratio=169]{beamer}

\mode<presentation>
\usetheme{Boadilla}
\definecolor{NTHU1}{RGB}{127,16,132}
\definecolor{NTHU2}{RGB}{228,210,231}
\definecolor{blue}{RGB}{30,90,205}
\definecolor{red}{RGB}{213,94,0}
\definecolor{green}{RGB}{0,128,0}
\definecolor{charcoalgray}{RGB}{108,104,104}
\setbeamercolor{title}{fg=NTHU1}
\setbeamercolor{frametitle}{fg=NTHU1}
\setbeamercolor{block title}{bg=NTHU1, fg=white}
\setbeamercolor{block body}{bg=NTHU2}
\setbeamercolor{structure}{fg=NTHU1}
\setbeamercolor{item projected}{fg=white}
\setbeamercolor{item}{fg=NTHU1}
\setbeamercolor{subitem}{fg=NTHU1}
\setbeamercolor{section in toc}{fg=NTHU1}
\setbeamercolor{description item}{fg=NTHU1}
\setbeamercolor{caption name}{fg=NTHU1}
\setbeamercolor{button}{bg=NTHU1, fg=white}
\setbeamercolor{caption name}{fg=NTHU1}
\usepackage{graphics}
\usepackage{tikz}
\usepackage{amsmath}
\usepackage{bbm}
\usetikzlibrary{decorations.pathreplacing}
\usepackage{geometry}
\usepackage{booktabs}
\usepackage{bookmark}
\usepackage{multirow, makecell}
\usepackage{float}
\usepackage{fancyvrb}
\usepackage{caption}
\captionsetup[figure]{singlelinecheck = false, format= hang, justification=centering}
\captionsetup[subfigure]{singlelinecheck = false, format= hang, justification=raggedright}
\usepackage{subcaption}
\usepackage{adjustbox}
\usepackage{hyperref}
\usepackage{threeparttable}
\usepackage{appendixnumberbeamer} 
\usepackage[T1]{fontenc}
\usepackage[default]{lato} 
\usepackage{smartdiagram}
\smartdiagramset{
  back arrow disabled = true,
  module minimum width=1cm,
  module x sep = 3,
  uniform arrow color=true,
}
\usepackage{xcolor}
\usepackage{colortbl}
	\definecolor{light-gray}{gray}{0.99}

\newenvironment{wideitemize}{\itemize\addtolength{\itemsep}{10pt}}{\enditemize}
\newenvironment{wideenumerate}{\enumerate\addtolength{\itemsep}{10pt}}{\endenumerate}
\newenvironment{widedescription}{\description\addtolength{\itemsep}{10pt}}{\enddescription}
\hypersetup{
colorlinks=true,
linkcolor=NTHU1,
filecolor=green, 
urlcolor=blue,
}
\beamertemplatenavigationsymbolsempty
\setbeamercolor{author in head/foot}{bg=white, fg=NTHU1}
\setbeamercolor{title in head/foot}{bg=white, fg=NTHU1}
\setbeamercolor{date in head/foot}{bg=white, fg=NTHU1}
\setbeamercolor{section in head/foot}{bg=white, fg=NTHU1}
\setbeamercolor{headline}{bg=white}
\setbeamertemplate{footline}{
    \leavevmode%
    \hbox{%
        \begin{beamercolorbox}[wd=.333333\paperwidth,ht=2.25ex,dp=1ex,center]{date in head/foot}%
            \usebeamerfont{date in head/foot}%\insertdate
        \end{beamercolorbox}%
        \begin{beamercolorbox}[wd=.333333\paperwidth,ht=2.25ex,dp=1ex,center]{title in head/foot}%
            \usebeamerfont{title in head/foot}\insertshorttitle
        \end{beamercolorbox}%
        \begin{beamercolorbox}[wd=.333333\paperwidth,ht=2.25ex,dp=1ex,right]{date in head/foot}%
            \usebeamerfont{date in head/foot}\insertframenumber{}\hspace*{2em}
        \end{beamercolorbox}}%
        \vskip0pt%
    }
    %% May need to script the introduction!!!!!!
%\setbeamercolor{page number in head/foot}{fg=NTHU1}
\setbeamertemplate{section in toc}[sections numbered]
\setbeamertemplate{subsection in toc}{\leavevmode\leftskip=3em\rlap{\hskip-1.75em\inserttocsectionnumber.\inserttocsubsectionnumber}
\inserttocsubsection\par}
\setbeamerfont{subsection in toc}{size=\footnotesize}
%\setbeamertemplate{headline}{%
  %\begin{beamercolorbox}[ht=5.5ex]{section in head/foot}
    %\vskip2pt\insertnavigation{0.33\paperwidth}\vskip2pt
  %\end{beamercolorbox}%
%}
\newenvironment{transitionframe}{\setbeamercolor{background canvas}{bg=white}\setbeamertemplate{footline}{} \begin{frame}[noframenumbering]}{\end{frame}}


\makeatletter
\let\@@magyar@captionfix\relax
\makeatother


\title[]{Conflict and Peace} % Change this regularly
\subtitle[]{Week 1: Introduction to the course and empirical analysis}
\author[Lee]{Seung-hun Lee }
\institute[]{Taipei School of Economics\\ National Tsing Hua University}

\date[]{September 5th, 2025}

\begin{document}
\begin{transitionframe}
\titlepage
\end{transitionframe}


%%% Color slides for section headers: Use for colloquium version (The ones bewteen \iffals and \fi)

\section{Introduction}
\begin{transitionframe}
\frametitle{Roadmap for today}
\setlength{\parskip}{1.25ex}
\tableofcontents[currentsection, hideothersubsections]
\end{transitionframe}

\subsection{Course logistics}
\begin{frame}
\frametitle{Course logistics} 
\begin{wideitemize}
\item Instructor: Seung-hun Lee 
\begin{itemize}
\item Office location: Room C06 in this building
\item Office hours: Thursdays 3:30-5:30PM \href{https://www.calendly.com/seunghunlee0918}{(Sign up here)}
\item From South Korea
\item Ph.D. in Economics from Columbia University (2024)
\end{itemize}
\item TA: Su Mon Htet Naing
\begin{itemize}
\item Contact: ekaterina.naing@gmail.com
\end{itemize}
\item Class time and location
\begin{itemize}
\item Time: Fridays 9AM - 12PM
\item Location: Room A09
\end{itemize}
\item Coming prepared to the lecture
\begin{itemize}
\item Lecture notes are uploaded a day in advance (No separate textbooks)
\item Keep up with the readings: more on that later
\end{itemize}
\end{wideitemize}
\end{frame}

\subsection{Topics covered in course}
\begin{frame}
\frametitle{Course goal: Viewing politics from an economic perspective} 
\begin{wideitemize}
\item Focus on political economy of development related to conflict and peace
\begin{itemize}
\item Conflicts: Violent events that disrupt economic and social stability - crime, political violence, civil wars, inter-state wars
\item Peace: Absence of conflict, or development of economic and institutional arrangements that prevents conflict
\end{itemize}
\item Economic perspective: Making causal inferences between events/outcomes
\begin{itemize}
\item Understand the theories behind incentives of stakeholders involved
\item Able to make educated, empirical analysis that identifies causal relations
\end{itemize}
\item Goals and requirements
\begin{itemize}
\item Be able to critically engage with scholarly work, conduct independent empirical analysis, and develop original research
\item Familiarity with statistical inference is necessary!
\item \textit{Helpful: Graduate-level micro/macroeconomics}
\end{itemize}
\end{wideitemize}
\end{frame}


\begin{frame}
\frametitle{Course topics (Refer to syllabus for full reading list)} 
\begin{wideitemize}
  \item Course introduction and empirical analysis tools (Week 1)
\item Cause of conflicts (Week 2-4)
\begin{itemize}
\item \textit{How do economic incentives drive individuals and states to conflict?}
\item \textit{How do political and ethnic dynamics lead to conflict?}
\item \textit{How and why do people engage in criminal activites?}
\end{itemize}
\item Consequences of conflicts (Week 5-7)
\begin{itemize}
\item \textit{How do they affect individuals?}
\item \textit{How do they shape economic development?}
\item \textit{How do they alter the institutional development of the state?}
\end{itemize}

\end{wideitemize}
\end{frame}

\begin{frame}
\frametitle{Course topics (Refer to syllabus for full reading list)} 
\begin{wideitemize}

\item Responses to conflicts (Week 9-10)
\begin{itemize}
\item \textit{How do individuals overcome collective action problems?}
\item \textit{What drives people to migrate?}
\item \textit{Hoow do people integrate to the new society?}
\end{itemize}
\item State capacity (Week 11-13)
\begin{itemize}
\item \textit{Why do some states fail?}
\item \textit{How do capable states build and maintain fiscal capacity?}
\item \textit{How do capable states manage personnel capacity?}
\end{itemize}

\end{wideitemize}
\end{frame}




\subsection{Course grading, tasks, and timeline}
\begin{frame}
\frametitle{Tasks and grading scale} 
\begin{wideitemize}
\item Weekly snap presentation (10\% of the grades)
\begin{itemize}
\item Two students will give a 15-minute presentation on two papers from the reading list marked with double asterisks
\item The presenters will be chosen at random on the day of the class - \textbf{So everyone must prepare the slides before the class}
\item Use \textbf{no more} than 8 slides, upload slides to eeclass after class
\end{itemize}
\item Replication exercise (20\% of the grades) - Due on Nov 21st
\begin{itemize}
\item You are asked to use R or Stata to replicate Brodeur (2018), or Campante, Do (2014) - further details on these papers in syllabus
\item You will produce key results (in tables and figures) from one of those papers, and submit a \textbf{replication report} and \textbf{codes}
\item Report whether the data is accessible, and the key findings replicate. 
\item More important: \textbf{Make your own extension} - alternative regression specifications, different independent variables, different falsification tests
\end{itemize}
\end{wideitemize}
\end{frame}

\begin{frame}
\frametitle{Tasks and grading scale} 
\begin{wideitemize}
\item Virtual paper (50\% of the grades)
\begin{itemize}
\item Develop a viable paper plan that clearly motivates the research question, explains the background and data, empirical strategy, and preliminary findings
\item Does not need to be a complete paper - \textbf{A half-finished but original and exciting project is preferred} over a complete yet uninteresting one
\item Start with a two-page proposal (10\% of the grades) - Due October 24 
\item Then, 30-minute presentation on the full paper (10\% of the grades) - December 5th 
\item Based on comments, finish the draft (30\% of the grades) - December 21st
\end{itemize}
\item Final Exam (20\% of the grades) - December 19th 
\begin{itemize}
\item Cumulative exam that will cover papers in single/double asterisks
\item Exact format to be determined soon
\end{itemize}
\end{wideitemize}
\end{frame}

\begin{frame}
\frametitle{More on the virtual paper} 
\begin{wideitemize}
\item Format of the virtual paper: Mimics many scholarly articles
\begin{enumerate}
\item \underline{Introduction}: What is your research question? Why is it important from a scholarly or policy standpoint?
\item \underline{Literature review}: What is the gap in the literature that you relate your work to? How are you filling that gap?
\item \underline{Background}: What is the relevant information on the context that you are studying?
\item \underline{Data}: What data source are you using or planning to use? Why should it be that dataset? 
\item \underline{Empirical strategy}: What statistical methods do you use? How do they allow for causal inference? What are the potential threats to that inference?
\item \underline{Preliminary results}: What are the statistical relationships that you have found so far?
\item \underline{Mechanism and falsification tests}: What factors may explain your main results? What are your plans to address potential factors that may nullify your findings?
\end{enumerate}
\end{wideitemize}
\end{frame}

\begin{frame}
\frametitle{Important deadlines} 
\begin{wideitemize}
\item October 24th: Presentation on 2-page proposal (submit draft before class)
\item November 21st: Replication exercise (submit report and codes by email)
\item December 5th: Virtual paper presentation
\item (Week of) December 19th: Final Exam (exact time TBA)
\item December 21st: Submit draft of the virtual paper by midnight
\end{wideitemize}
\end{frame}




\subsection{Rules and Tips}
\begin{frame}
\frametitle{Some ground rules} 
\begin{wideitemize}
\item Usage of AI: Strictly for assistive purposes
\begin{itemize}
\item Allowed for proofreading, coding support
\item Do not generate complete assignments using AI 
\item All writings, especially virtual paper, must be original
\item To show that you are working on your own idea, regularly discuss your ideas with me through office hours
\end{itemize}
\item Collaboration: You can collaborate, under these two conditions
\begin{itemize}
\item Each student must write up their own assignments
\item Students must clearly identify their collaborators
\end{itemize}
\end{wideitemize}
\end{frame}





% Second paragraph of intro
\section{Empirical Analysis}
\begin{transitionframe}
\setlength{\parskip}{1.25ex}
\tableofcontents[currentsection, hideothersubsections]
\end{transitionframe}



\subsection{Conceptualizing Treatment Effects}
\begin{frame}
\frametitle{Causal inference in social sciences: Quantifying causal relationships} 
\begin{wideitemize}
\item In this class (and beyond), we aim to clarify quantitative, causal relationship between different events and outcomes
\begin{itemize}
\item \textit{Does higher commodity prices reduce wars? (Caselli et al 2015; Dube, Vargas 2013)}
\item \textit{Does extractive colonial institutional backdrop hurt current institutional development? (Acemoglu et al 2001; Michalopoulos, Papaiaonnou 2016)}
\item \textit{Does negative shocks to economic productivity lead to revolutions (Brückner, Ciccone 2011)}
\end{itemize}
\item Treatment effect measures the size of impact these events have on sample of interest
\begin{itemize}
\item Identify an event that induces some shocks to individual behavior/societal outcome
\item Find a suitably comparable sample (i.e. the control group and treatment)
\item Compare outcomes for those affected by the event to those not affected 
\end{itemize}
\item To that end, we collect data from a suitably defined population and use various methods to estimate the effects of these events.
\end{wideitemize}
\end{frame}

\begin{frame}
\frametitle{Correlation vs Causation}
\begin{wideitemize}
\item Suppose you have random variables $X$ and $Y$. You want to identify if $X$ \textit{causes} $Y$
\item A \textbf{correlation} between $X$ and $Y$ is a statistical measure that describes how the two variables move together
\begin{itemize}
\item It captures \textit{any} type of statistical dependence that moves the two variables together: Causation, but also others too!
\item Not causal 1: $Y$ can cause $X$ (reverse causality)
\item Not causal 2: $X$ and $Y$ co-move because of $Z$ that affects both (Omitted variable bias)
\end{itemize}
\item A \textbf{causal} relationship: \textit{Cleanly (exogenously)} changing variables $X$ leads to changes in $Y$
\begin{itemize}
\item Changes in $X$ may be a combination of changes from $X$ alone + other factors that co-move with $X$ and affect $Y$ (You do not want the latter)
\item Randomized experiments: One way of Isolating clean changes in $X$ 
\item Econometrics: Concisely explains the relationship between $X$ and $Y$ variables in a single equation (Even when experiment is not feasible)
\end{itemize}
 \end{wideitemize}
\end{frame}

\begin{frame}
\frametitle{Treatment effect in mathematics} 
\begin{wideitemize}
\item Mathematical notations
\begin{itemize}
  \item $Y_i$: observed outcome for individual $i$ % e.g. Income, votes, vote particpation, health
  \item $D_i=d$: Treatment assignment ($d=1$ for treated, 0 for control group) 
  \item $Y_i(d)$: Outcome for treated individual if $d=1$, control group if $d=0$. 
\end{itemize}
\item Treatment effect: How much does outcome differ between treated and control group
\[
  TE = Y_i(1)-Y_i(0)
\]
\begin{itemize}
\item With large number of $i$, we want an average of TE (henceforth average treatment effect)
\[
ATE = E[Y_i(1)-Y_i(0)]=E[Y_i(1)]-E[Y_i(0)]
\]
\item Fundamental problem: If $i$ is treated, it cannot be in the control group 
\item So either one of $Y_i(1)$ or $Y_i(0)$ is missing! $\to$ We do not exactly know the \textbf{counterfactual}
\item Econometrics provides ways to provide sound judgement on what the counterfactual looks like  
\end{itemize}
\end{wideitemize}
\end{frame}



\begin{frame}
\frametitle{Potential outcomes and randomized experiments} 
\begin{wideitemize}
\item Potential outcomes provides framework to match observations to treatment effect
\[ 
\begin{aligned}
Y_ i &= Y_i(1)D_i + Y_i(0)(1-D_i)\\
&=Y_i(0)+D_i(Y_i(1)-Y_i(0))\\
\end{aligned}
\]
\item In terms of conditional expectations, 
\[ 
\begin{aligned}
E[Y_ i|D_i] &=E[Y_i(0)+D_i(Y_i(1)-Y_i(0)) |D_i]   \\
&=E[Y_i(0)|D_i]+E[Y_i(1)-Y_i(0)|D_i]D_i \\
\end{aligned}
\]
\item Note that 
\begin{itemize}
\item If $D_i=1$, then $Y_i = Y_i(1)$ and $E[Y_i|D_i=1]=E[Y_i(1)|D_i=1]$
\item If $D_i=0$, then $Y_i = Y_i(0)$ and $E[Y_i|D_i=0]=E[Y_i(0)|D_i=0]$
\end{itemize}
\end{wideitemize}
\end{frame}



\begin{frame}
\frametitle{Potential outcomes and randomized experiments} 
\begin{wideitemize}

\item With randomized experiments, we assume that the counterfactual is identical \textbf{except} for treatment status
\begin{itemize}
\item Mathematically, treatment assignment is independent of potential outcomes
\[
E[Y_i(d)]=E[Y_i(d)|D_i=1]=E[Y_i(d)|D_i=0] \ \text{for } d\in\{0,1\}
\]
\item This implies that 
\[
\begin{aligned}
ATE &= E[Y_i(1)]-E[Y_i(0)]\\
&=E[Y_i(1)|D_i=1]-E[Y_i(0)|D_i=0]\\
&=E[Y_i|D_i=1]-E[Y_i|D_i=0] (\because\text{previous slide})\\
\end{aligned}
\]
or that ATE is obtained by subtracting average value of $Y_i$ for treatment group by that of control group (which we can do with the observed data we have)
\end{itemize}
\end{wideitemize}
\end{frame}

\subsection{Basic Regressions: Ordinary Least Squares}
\begin{frame}
\frametitle{Ordinary Least Square (OLS): Mathematical approach } 
\begin{wideitemize}
\item Write 
\[
Y_i = \alpha + \beta X_i + u_i 
\] 
where $X_i$ is an independent variable and $u_i$ is the error (determinants of $Y_i$ not captured by $X_i$) 
\item Exogeneity: We assume that $u_i$ is uncorrelated with $X_i$, $E[u_i|X_i]=0$ (very important!)
\item OLS estimates the parameter of interest $\beta$ by minimizing the sum squared of errors, giving us the $\beta$ that best matches the patterns in the observation
\[
\hat{\beta}_{OLS} = \min_\beta \sum_{i=1}^n (Y_i - \alpha-\beta X_i)^2
\]
\begin{itemize}
\item We can do the same for $\alpha$, which capture baseline average (average when all independent variables are zero)
\end{itemize}
\end{wideitemize}
\end{frame}

\begin{frame}
\frametitle{Conceptualizing OLS} 
\begin{figure}
\includegraphics[width=0.67\textwidth, keepaspectratio]{ols.png}
\end{figure}
\end{frame}




\begin{frame}
\frametitle{Ordinary Least Square (OLS): Potential outcomes approach } 
\begin{wideitemize}
\item Recall the potential outcome framework 
\[
\begin{aligned}
Y_ i &=Y_i(0)+D_i(Y_i(1)-Y_i(0))\\
E[Y_ i|D_i]&=E[Y_i(0)|D_i]+E[Y_i(1)-Y_i(0)|D_i]D_i\\
\end{aligned}
\] 

\item Combined with OLS regression
\[
\begin{aligned}
Y_i &= \alpha + \beta D_i + u_i \\
\implies E[Y_i|D_i]&=\alpha+\beta D_i+E[u_i|D_i] \\
&=\alpha+\beta D_i\\
\end{aligned}
\] 
\begin{itemize}
\item With random assignment, $ATE=E[Y_i(1)-Y_i(0)|D_i]=E[Y_i(1)-Y_i(0)]=\beta$
\item Thus, ATE can be obtained with OLS estimates of $\beta$, or $\hat{\beta}_{OLS}$
\end{itemize}
\item You can add more variables, include $X_i$ on top of $D_i$ for better precision
\end{wideitemize}
\end{frame}


\begin{frame}
\frametitle{Internal and external validity problems: Why OLS is not always usable} 
\begin{wideitemize}
\item \textbf{External validity} is concerned with the applicability to other contexts.  \medskip
\item \textbf{Internal validity} assesses whether the model accurately represents the sample
\begin{itemize}
\item \textbf{Omitted variable bias: } Did the researcher left out factors that could affect $Y$ and are correlated with $X$?  $\to$ \textcolor{blue}{Multivariate regression, panel regression}
\item \textbf{Wrong functional form: } Did the researcher incorporate $X$ in a proper functional form? (quadratic, logs?) $\to$ \textcolor{blue}{add logs, or use logit/probit}
\item \textbf{Measurement error: } If $X$ does not have an accurate measure, OLS results in attenuation bias   $\to$ \textcolor{blue}{Instrumental variable}
\item \textbf{Sample selection bias: }  If certain groups are more likely to not respond, then the data fails to estimate the true effects $\to$ \textcolor{blue}{Instrumental variable, Natural experiments} 
\item \textbf{Simultaneous causality bias: } There are cases where $Y$ causes $X$ too, also leading to the bias (consider supply and demand) $\to$ \textcolor{blue}{Instrumental variable}
\end{itemize}
\item Controlled experiment nearly impossible in political economy $\to$ Natural experiment that uses external events/shock as treatment
\end{wideitemize}
\end{frame}


\subsection{Panel Regressions}
\begin{frame}
\frametitle{When exogeneity fails} 
\begin{wideitemize}
\item  $E[u_i|D_i]=0$ rarely holds: Treatment is not always random, or some individuals are more likely to be treated than others
\begin{itemize}
\item This biases our results: Suppose that $E[u_i|D_i]=\phi D_i + v_i$ ($u_i$ and $D_i$ have correlation)
 \[
 \begin{aligned}
  E[Y_i|D_i]&=\alpha+\beta D_i+E[u_i|D_i]\\
  &=\alpha + (\beta+\phi)D_i + v_i\\
 \end{aligned}
\]
With OLS on $Y_i$ and $D_i$, $\hat{\beta}$ captures $\beta+\phi$ - Treatment effect and \textcolor{red}{endogeneity bias}
\end{itemize}
\item When and why do biases happen?
\begin{itemize}
\item Treated groups may be a good performer to begin with (and vice versa)
\item Differences in baseline characteristics besides treatment $\to$ inaccurate TE estimates
\end{itemize}
\item Now what? Control for $\phi$ with individual/time indicators (Panel), exogenous proxy (IV), or natural experiments
\end{wideitemize}
\end{frame}






\begin{frame}
\frametitle{Panel regression fixed effects approach} 
\begin{wideitemize}
\item We observe multiple individuals for multiple time periods ($i\in \{1,...,N\}, t\in \{1,...,T\}$)
\[
Y_{it} = \beta_0 + \beta_1X_{1,it}+ ... +\beta_kX_{k,it}+u_{it}
\]
\item Problem: Unobserved individual characteristics and time trends could affect outcomes and cause omitted variable bias
\begin{itemize}
\item Individual's unobserved ability/characteristics (e.g. culture)
\item A bad shock in year $t$ that affects everyone (COVID years)
\end{itemize}
\item Panel regression allows us to capture these using individual and time fixed effects\
\[
Y_{it}=\beta_1 X_{1,it}+...+\beta_k X_{k,it} + \alpha_i +\lambda_t + u_{it}
\]
\begin{itemize}
\item $\alpha_i$: Dummy for each individual $i\in\{1,...,N\}$
\item $\delta_t$: Dummy for each time $t\in\{1,...,T\}$
\end{itemize}
\end{wideitemize}
\end{frame}

\subsection{Instrumental Variable Regressions}
\begin{frame}
\frametitle{Instrumental variables: Exogeneity amidst endogeneity} 
\begin{wideitemize}
\item OLS hinges on independent variable $X_i$ being uncorrelated with $u_i$
\begin{itemize}
\item This often breaks: omitted variable bias, measurement error, and reverse causality etc..
\item Parts of $X_i$ that are \textcolor{red}{correlated with $u_i$}: What we want is part of $X_i$ \textcolor{blue}{independent of $u_i$}.
\end{itemize}
\item Solution is to find another variable $Z_i$ that are
\begin{itemize}
\item (1) Correlated with $X$ but (2) not with $u$: This gets us part of $X_i$ \textcolor{blue}{independent of $u_i$}
\item (1) is the relevancy condition, (2) is the exogeneity condition
\item Challenge: Exogeneity is difficult to argue (usually argue that $Z_i$ affects $Y_i$ only through $X_i$)
\end{itemize}
\item This is estimated with two stage least squares (2SLS)
\begin{itemize}
\item Original Equation: $Y_i = \alpha + \beta X_i + u_i$
\item First stage: $X_i=\delta_0 + \delta_1Z_i+v_i \to$ Get $\widehat{X}_i=\hat{\delta}_0+\hat{\delta}_1Z_i$ ($\widehat{X}_i$ is uncorrelated with $u_i$ )
\item Second stage: $Y_i = \beta_0 + \beta_1 \widehat{X}_i + u_i$
\end{itemize}
\end{wideitemize}
\end{frame}

\begin{frame}
\frametitle{Instrumental variables using natural environment} 
\begin{wideitemize}
\item Source: Brückner, Ciccone, ``Rain and the Democratic Window of Opportunity'', \textit{Econometrica}, 79(3), 923-947

\item Question: Can transitory shocks to income ($X_i$) lead to democratic changes ($Y_i$)?

\item Income shocks can be endogenous: Places with higher income are likely more democratic to begin with
\begin{itemize}
\item Use (lack of) rainfall ($Z_i$): Correlates with income through agricultural productivity but unlikely to affect democratic transitions alone
\item Plausible: Very hard for rainfall itself to lead to democratization
\end{itemize}

\item Result: 1\% transitory negative income shock $\to$ 1.3 percentage point rise in democratic transitions

\end{wideitemize}
\end{frame}

\begin{frame}
  \centering
\frametitle{Instrumental variables using natural environment} 
\includegraphics[width=0.75\textwidth]{iv.jpeg}
\end{frame}



\subsection{Natural Experiments: Difference-in-Differences, Event studies, and Regression Discontinuity}
\begin{frame}
\frametitle{Natural experiments: Difference-in-differences} 
\begin{wideitemize}
\item Assume two time periods (before and after) and two groups (treated and controlled)
\item You have the following regression
\[
Y_{it} = \beta_0 + \beta_1 D_{it} + \beta_2 P_{it} + \beta_3 D_{it}P_{it}+u_{it}
\]
where $D=1$ if in treatment group, $P=1$ if after treatment started
\item You can get the following group averages
\begin{itemize}
\item Treated, After: $E[Y_{it}|D_{it}=1, P_{it}=1] = \beta_0+\beta_1 + \beta_2 + \beta_3$
\item Treated, Before: $E[Y_{it}|D_{it}=1, P_{it}=0] =\beta_0+ \beta_1 $ 
\item Controlled, After: $E[Y_{it}|D_{it}=0, P_{it}=1] =\beta_0+ \beta_2 $
\item Controlled, Before: $E[Y_{it}|D_{it}=0, P_{it}=1] =\beta_0 $ 
\end{itemize}
\item DID: $[\bar{Y}_{\text{treat,after}}-\bar{Y}_{\text{treat,before}}]-[\bar{Y}_{\text{control,after}}-\bar{Y}_{\text{control,before}}]  \to \hat{\beta}_3$
\item Can be generalized to many time periods and many groups with TWFE, Event-studies etc (active area of research)
\end{wideitemize}
\end{frame}

\begin{frame}
\frametitle{Natural experiments: Conceptualizing Difference-in-differences} 
\begin{figure}
\includegraphics[width=0.7\textwidth, keepaspectratio]{did.png}
\end{figure}
\end{frame}

\begin{frame}
\frametitle{Natural experiments: Difference-in-Differences} 
\begin{wideitemize}
\item Source: Andreas V. Jensen (2025) ``Educating for Democracy? Going to College Increases Political Participation'', \textit{British Journal of Political Science}, 55,1-12
\item Question: Does going to college increase political participation?
\item Method: Uses a survey of US students who graduated high school at 2004, observe differences in election participation in 2004 and 2008
\[
\text{Turnout}_{it}=\alpha_i+\delta_t + \beta[\text{College}_i\times\text{Post}_t]+\gamma X_{it}+u_{it}
\]
\begin{itemize}
  \item Here $\alpha_i$ and $\delta_t$ are individual and time fixed effects\\ (Can you find their equivalent variables in the previous equation?)
  \item Treated: Those who entered college between 04-08
  \item Control: Did not go to college
\end{itemize}
\item Result: College enrollment increase particpation by 13 percentage points
\end{wideitemize}
\end{frame}

\begin{frame}
\frametitle{Natural experiments: Difference-in-Differences} 
\begin{figure}
\includegraphics[width=0.7\textwidth, keepaspectratio]{did_table.png}
\end{figure}
\end{frame}

\begin{frame}
\frametitle{Natural experiments: Regression discontinuity} 
\begin{wideitemize}
\item Comparing individuals just entering the treatment vs marginally did not receive one
\begin{itemize}
\item Assume you are treated if $Z_i \geq c$ and not if otherwise
\item Then, you run this regression on $Z_i \in [c-h, c+h]$ for some small $h$
\[
Y_i = \beta X_i + \gamma (1[Z_i\geq c]\times X_i) + \delta (X_i \times (Z_i-c)) + \mu (X_i \times (Z_i-c)\times 1[Z_i\geq c])+u_i
\]
For simplicity, assume that no other independent variables are in place: $X_i=1$
\item $1[Z_i \geq c]=1$ if $i$ passes threshold for eligibility (0 if otherwise)
\item $\gamma$ captures the treatment effect: Difference between those in and out of treatment
\item We also allow slope to be different under and above the cutoff: $\delta$ for people below cutoff and $\delta+\mu$ for those above
\end{itemize}
\item Treatment effect estimated here is local: May not hold for out-sample
\item There may be cases where threshold is not strict (fuzzy RD)
\end{wideitemize}
\end{frame}


\begin{frame}
\frametitle{Regression discontinuity in use with electoral rules} 
\begin{wideitemize}
\item Source: Thomas Fujiwara (2011) ``A Regression Discontinuity Test of Strategic Voting and Duverger’s Law'', \textit{Quarterly Journal of Political Science}, 6, 197-233
\item Question: Do third-party candidates perform well in majoritarian elections with run-offs than ones with no run-offs?
\item Design: In Brazil, precincts with more than 200,000 registered voters have run-off elections for mayors (otherwise, no run-offs)
\item Findings: Third-party candidates perform signficantly worse in places without run-offs, suggesting that voters consider winning probabilities of candidates into account
\end{wideitemize}
\end{frame}


\begin{frame}
\frametitle{Natural experiments: Regression discontinuity} 
\begin{figure}
\includegraphics[width=0.8\textwidth, keepaspectratio]{rd.png}
\end{figure}
\end{frame}


%%%%%%%%%%%
\end{document}
